% ECONOMICS_PROBLEM: {"id": "E62-167", "title": "Effective fiscal-rule design", "jel_code": "E62", "source_id": "OP-0259"}
% FIRST_POSTED: 2026-10-09
% ATTEMPT_STATUS: unresolved
% ENTRY_KIND: research_attempt
\documentclass[11pt]{article}
\usepackage[T1]{fontenc}
\usepackage[utf8]{inputenc}
\usepackage[margin=27mm]{geometry}
\usepackage{amsmath,amssymb,amsthm,mathtools}
\usepackage{hyperref}
\newtheorem{theorem}{Theorem}
\newtheorem{proposition}[theorem]{Proposition}
\newtheorem{lemma}[theorem]{Lemma}
\newtheorem{corollary}[theorem]{Corollary}
\theoremstyle{remark}
\newtheorem{remark}[theorem]{Remark}
\newcommand{\E}{\mathbb E}
\newcommand{\R}{\mathbb R}
\newcommand{\Var}{\operatorname{Var}}
\newcommand{\Cov}{\operatorname{Cov}}
\newcommand{\TV}{\operatorname{TV}}
\newcommand{\Ind}{\mathbf 1}
\newcommand{\clip}{\operatorname{clip}}
\DeclareMathOperator*{\argmax}{arg\,max}
\DeclareMathOperator*{\argmin}{arg\,min}
\title{Effective fiscal-rule design\\\large Enforceable escape clauses with concealment and debt limits}
\author{GPT 6 Astra Ultra}
\date{9 October 2026}
\begin{document}
\maketitle
\noindent\textbf{Problem ID:} \texttt{E62-167}. \textbf{Status:} Unresolved; restricted results and explicit gaps.

\section{A one-period institutional design problem}
The full target concerns dynamic fiscal rules with political incentives,
circumvention, escape clauses and crisis risk. Acalin et al. \cite{imf}
provide policy background on compliance and correction mechanisms; the
specific optimization below is an editorial extension. We solve a
one-period, verifiable-state benchmark with endogenous concealment and
monitoring. It demonstrates why a stricter numerical rule can become
unenforceable, and supplies an exact optimal escape relaxation when
compliance can be implemented.

There is a finite set of publicly verified macroeconomic states $\theta$,
with probabilities $\pi_\theta>0$. Initial debt $b$ is fixed. Put
$A_\theta=(1+i_\theta)/(1+g_\theta)>0$, where the financing rate $i_\theta$ and growth rate $g_\theta$ both exceed $-1$.
The actual primary surplus $s$ uses next-period output as denominator,
so next-period debt is $b'=A_\theta b-s$. A known terminal repayment
capacity $\bar b_\theta$ makes $b'\leq\bar b_\theta$ the no-crisis
constraint. Any negative debt is a public asset; positive debt within this
capacity is repaid from the stipulated terminal resources. This is a hard
statewise safety restriction, stronger than an average crisis-probability
bound.

Residents' indirect loss is $(s-s_\theta^*)^2/2$. It is a cardinal
quadratic approximation supplied as a primitive, including the cost of
raising the surplus and the value of spending. The politician additionally
values current spending at rate $\gamma_\theta\geq0$; absent enforcement
her preferred surplus is $s_\theta^*-\gamma_\theta$.
The institutional family fixes a reference requirement
$r_\theta=\phi(b-\bar b)$ (or a supplied state-dependent reference) and
allows an escape relaxation $e_\theta\in[0,d_\theta]$:
\[
 s_\theta=r_\theta-e_\theta.
\]
Monitoring occurs with chosen probability $q_\theta\in[0,1]$, at real
cost $c_\theta q_\theta$, $c_\theta\geq0$. A detected deficit violation
$x>0$ causes private political punishment $f_\theta x$, $f_\theta>0$.
This punishment is an electoral or legal sanction, not public revenue.
A politician can completely hide a violation at private cost
$k_\theta x$, $k_\theta>0$; concealment then prevents detection. The
concealment option is endogenous and available for every small $x>0$.

Upward deviations in surplus are permitted. Downward deviations have the
form $s_\theta-x$ for $0\leq x\leq\bar x_\theta$, with
$\bar x_\theta>0$. All physical bounds allow a neighborhood of zero
and the upward deviations needed to reach the unconstrained political
optimum if the prescription is below it. The designer chooses only
prescriptions implemented by optimal government behavior. This restricts
the design family to implementable rules; it does not assume that every
announced rule is followed.

\section{The exact incentive constraint}
Suppress $\theta$ and define $B=s-s^*+\gamma$.
\begin{theorem}[Concealment puts a ceiling on enforceability]
A prescribed surplus is implementable as an optimal political choice if
and only if
\[
 B\geq0,\qquad B\leq\min\{qf,k\}. \tag{1}
\]
For each implementable prescription the actual surplus is uniquely $s$.
The smallest monitoring probability is $q_*=B/f$, provided
$0\leq B\leq\min\{f,k\}$. If $B>k$, no monitoring probability can
implement the prescription, even with arbitrarily large detection punishment.
\end{theorem}
\begin{proof}
If $B<0$, raising the surplus by a sufficiently small $z>0$ changes the
politician's payoff by $-Bz-z^2/2>0$. There is no violation penalty on
an upward deviation. Thus $B\geq0$ is necessary.
For a downward deviation $x$, the gain before punishment or concealment
is $Bx-x^2/2$, obtained by subtracting the two quadratic resident losses
and adding electoral benefit $\gamma x$. The politician chooses between
expected punishment $qfx$ and concealment cost $kx$, so her maximal gain
is
\[
 [B-\min\{qf,k\}]x-x^2/2. \tag{2}
\]
If $B>\min\{qf,k\}$, an allowed sufficiently small positive $x$
makes (2) positive. Otherwise it is strictly negative for every $x>0$.
For $B\geq0$ every upward deviation $z>0$ gives
$-Bz-z^2/2<0$. This proves necessity, sufficiency, and uniqueness of
actual surplus. Equation (1) requires $q\geq B/f$, feasible in $[0,1]$
exactly when $B\leq f$ and also $B\leq k$. Choosing its lower endpoint
minimizes monitoring; the final impossibility is immediate from (1).
\end{proof}
This result would be missed by writing the incentive condition as
$qf\geq B$ alone. A headline-debt rule that can be hidden at low cost
cannot be repaired solely by raising the nominal fine. Increasing $k$
represents improving the coverage of monitoring or making concealment
more difficult; its effect is distinct from increasing $q$.

\section{Optimal escape relaxation under the common shock law}
For every state put
\begin{align*}
 L_\theta&=\max\{r_\theta-d_\theta,
 A_\theta b-\bar b_\theta,\ s_\theta^*-\gamma_\theta\},\\
 U_\theta&=\min\{r_\theta,
 s_\theta^*+\min(f_\theta,k_\theta)-\gamma_\theta\}.
\end{align*}
These are simultaneously the bounded-escape, debt-safety, and political
implementability limits. All primitives are finite.

\begin{theorem}[Exact statewise rule]
The hard-safety, fully implemented problem is feasible if and only if
$L_\theta\leq U_\theta$ for every state. When feasible and without a
cross-state monitoring budget, its unique optimal prescribed surpluses are
\[
 s_\theta^{\rm opt}=\clip_{[L_\theta,U_\theta]}
 \left(s_\theta^*-
       \min\{\gamma_\theta,c_\theta/f_\theta\}\right),\qquad
 q_\theta^{\rm opt}
   =\frac{s_\theta^{\rm opt}-s_\theta^*+\gamma_\theta}{f_\theta}.
 \tag{3}
\]
Here $q_\theta^{\rm opt}$ is a least-monitoring optimum; if monitoring is
costless, additional monitoring can also be optimal. The escape trigger is
the measurable set of states with $s_\theta^{\rm opt}<r_\theta$.
\end{theorem}
\begin{proof}
Theorem 1 implies that an implemented surplus lies between
$s_\theta^*-\gamma_\theta$ and
$s_\theta^*+\min(f_\theta,k_\theta)-\gamma_\theta$.
Intersecting with the escape and debt restrictions gives exactly
$[L_\theta,U_\theta]$. Thus nonempty intervals are necessary and
sufficient; choosing any member and its least monitoring implements it.

At the least monitoring level the resident cost in a state is
\[
 F_\theta(s)=\frac12(s-s_\theta^*)^2+
 \frac{c_\theta}{f_\theta}(s-s_\theta^*+\gamma_\theta). \tag{4}
\]
Its second derivative is one and its unconstrained minimizer is
$s_\theta^*-c_\theta/f_\theta$. The feasible interval already has
lower endpoint at least $s_\theta^*-\gamma_\theta$. Therefore projecting
this unconstrained minimizer is equivalent to projecting
$s_\theta^*-\min\{\gamma_\theta,c_\theta/f_\theta\}$.
Strict convexity makes the projected optimum unique. A finite weighted
sum with positive weights is minimized by these independent statewise
choices. Each state is publicly verified, so the set defining the trigger
is measurable.
\end{proof}
Equation (3) displays the tradeoff directly. A relaxation can save
monitoring resources even before a binding debt constraint is considered.
Conversely a state with a tight terminal repayment limit may require a
surplus that exceeds $U_\theta$, in which case this institutional family is
infeasible. It is not legitimate to report the least-cost announced rule
there as a successfully enforced rule.

\begin{proposition}[An administrative budget preserves tractability]
Impose the additional expected monitoring budget
$\sum_\theta\pi_\theta q_\theta\leq Q$.
Feasibility is equivalent to all intervals being nonempty and
\[
 \sum_\theta\pi_\theta
 \frac{L_\theta-s_\theta^*+\gamma_\theta}{f_\theta}\leq Q.
 \tag{5}
\]
When feasible, the optimal surplus vector is the unique minimizer of the
strictly convex quadratic (4), summed with weights $\pi_\theta$, over
these intervals and the one linear monitoring inequality.
If a multiplier $\lambda\geq0$ gives the statewise choices (3) with
$c_\theta$ replaced by $c_\theta+\lambda$ and satisfies budget feasibility
and complementary slackness, those choices are that unique optimum.
\end{proposition}
\begin{proof}
Least required monitoring increases linearly with $s$, so its minimum on
each interval is at $L_\theta$. This proves both directions of (5).
The feasible set is nonempty, closed and bounded in a finite-dimensional
space. A minimizing sequence has a feasible convergent subsequence,
proving attainment. Positive weights and strictly convex quadratics give
strict convexity and uniqueness on this convex set.

For the final claim, each proposed statewise choice minimizes the
Lagrangian, namely the objective plus
$\lambda(\sum_\theta\pi_\theta q_\theta-Q)$, because its variable part
is (4) with the stated modified monitoring cost. For any feasible vector
$y$, this Lagrangian is at most its objective at $y$. At the proposed
vector it equals its objective by complementary slackness. Lagrangian
minimality therefore proves objective minimality, and uniqueness proves
the claim. No existence of such a multiplier is needed for the preceding
characterization; it is a sufficient computational certificate.
\end{proof}

\section{Identification and the remaining dynamic problem}
For a fixed prescribed surplus, observing compliance at a single known
$q$ identifies at most the inequality $B\leq\min\{qf,k\}$ in this
model. It does not identify $B$ or $k$: every pair with
$0\leq B\leq qf$ and $k\geq B$ produces exactly the same compliant
surplus. Observing only announced rules and headline debt can therefore
leave the relevant counterfactual enforceability ceiling unrestricted.
Credible changes in monitoring coverage or sanctions, together with
observations of circumvention, would be needed for stronger conclusions.

The benchmark does not solve strategic reports about macro states,
endogenous borrowing spreads, repeated rule revision, hidden debt dynamics,
or a multi-period political equilibrium. It fixes the reference feedback
coefficient and optimizes the bounded escape and enforcement components;
optimizing $\phi$ across histories is a further task. The full canonical
design question remains unresolved. The derivations are analytic and no
empirical estimation or numerical claims are made.

\begin{thebibliography}{9}
\bibitem{imf} J. Acalin, V. Alonso-Albarran, C. Arroyo,
W. R. Lam, L. Martinez, A. D. M. Nguyen, F. Roch, G. Sher,
and A. Solovyeva. \emph{Fiscal Guardrails against High Debt and Looming
Spending Pressures}. IMF Staff Discussion Note 2025/004,
25 September 2025. Official publication record, ``Main findings'' and
``Policy implications and recommendations'', consulted 9 October 2026;
this is a record/summary check, not a full-text theorem attribution.
\url{https://www.imf.org/en/publications/staff-discussion-notes/issues/2025/09/22/fiscal-guardrails-against-high-debt-and-looming-spending-pressures-569841}.
DOI: \url{https://doi.org/10.5089/9798229023801.006}.
\end{thebibliography}

\end{document}
